% Generated from retained local raw data by scripts/paper.py.
\section{Experimental evaluation}
\label{sec:evaluation}
We evaluated (i) the properties observed by the current implementation,
(ii) its behaviour on manually specified synthetic performances, and
(iii) the cost of its local analysis path. All performances were deterministic
synthetic MIDI, not pianist recordings or a user study. The earlier draft's
timings were not used as reference values or targets.

\paragraph{Setup and provenance.}
We compiled the seven unchanged RML specifications with the vendored compiler
and exercised the actual service and browser \texttt{toPerformance} export.
The service files were staged byte-for-byte in the existing Docker layout;
application code and specifications were not repaired. The accessible project
Git metadata did not identify a Cadenza commit, so SHA-256 manifests identify
the evaluated sources and inputs. The experiments ran on 14 October 2026,
on an Intel Core i7-7700HQ (2.80\,GHz; four cores/eight logical processors),
with 15.51\,GiB reported RAM and Linux 6.8.0-138-generic.
The isolated environment used Python 3.11.0rc1, Partitura 1.9.0,
NumPy 1.26.4, SciPy 1.17.1, lxml 6.1.3, mido 1.3.3 and SWI-Prolog 9.1.21.
Complete runtime versions, resource limits and dependency inventories are retained.
Numerical-library thread limits were one; CPU frequency and background load
were not controlled.

\paragraph{Behaviour.}
An original eight-note MusicXML exercise contains a two-note chord, dynamics,
staccato, tenuto and a pedal span. We froze manually reasoned per-specification
verdicts before execution, separately labelling musical errors under a declared
exercise rubric. Across 28 cases there were 0 deviations from these predictions,
including one predicted articulation runtime error. The baseline was accepted
by all seven metrics. Wrong pitch, out-of-range pitch, velocity and duration
cases were rejected as predicted. Chord spread of 29\,ms was accepted and
30/31\,ms rejected. At 120 beats per minute, the composed pedal-boundary
association accepted +25\,ms and rejected +26\,ms for this fixture, despite
the RML's nominal 0.25-beat tolerance. Missing and duplicate-note cases were
accepted by all metrics, exposing the absence of general score coverage and
one-to-one matching obligations. Syncopated pedal checked event order without
an elapsed-time bound: a 300\,ms delayed press passed that metric, whereas
pre-attack pedal failed. Missing release failed pedal-toggle but passed
syncopated-pedal. A final marked note caused articulation exit~2 because its
null next-attack gap reached compiled arithmetic. This defect was retained,
and error diagnostics are provided separately from acceptance and rejection.

\paragraph{Local analysis cost.}
The repository scores \emph{Great Fairy Fountain} and
\emph{Merry-Go-Round of Life} yielded 216 and 2,324 parsed notes, respectively.
The latter contains 2,418 pitched XML noteheads and 94 ties, explaining the
parsed count. The generated traces contain 464 and 4946
note-on, note-off and control-change events. Onsets use the implementation's
own timing map; this is not independent alignment validation. Durations,
velocities and pedal events follow fixed documented rules, with no randomness
or adjustment based on verdicts. MIDI encoding resolution was 0.5\,ms.
Synthetic timestamps already have the intended score-time origin; live
capture/playback timestamp-origin alignment was not tested.

We used a monotonic high-resolution clock around each actual
\texttt{analyse} call, including validation, temporary score creation,
MusicXML parsing, enrichment, JSONL serialisation and writing, fresh Prolog
startup, seven checks, verdict decoding and cleanup. Browser capture and
conversion, MIDI decoding, HTTP transport, rendering, Python import/startup
and ahead-of-time RML compilation were excluded. Each workload received
three warm-ups and 30 measured calls, sequentially in one Python process;
every call started a new Prolog process. Python libraries and OS caches could
remain warm. Delegating stage timers and warning capture add retained
instrumentation overhead. Total request durations were independently measured,
not reconstructed by summing stage medians.

\begin{table}[t]
\centering
\resizebox{\linewidth}{!}{% Generated by scripts/summarise.py; all durations are seconds.
\begin{tabular}{lrrrrrrrrr}
\hline
Workload & Notes & Events & Warm & Runs & Median & Q1 & Q3 & Min & Max \\
\hline
score-2324 & 2324 & 4946 & 3 & 30 & 5.788159 & 5.689245 & 5.820485 & 5.579121 & 6.109774 \\
score-216 & 216 & 464 & 3 & 30 & 0.703565 & 0.688775 & 0.726932 & 0.674853 & 0.776490 \\
\hline
\end{tabular}
}
\caption{Local analysis times in seconds, with linear-interpolation quartiles.
Every measured sample is retained; warm-ups are excluded from the statistics.}
\label{tab:local-cost}
\end{table}

The median totals were 0.703565\,s and 5.788159\,s
(Table~\ref{tab:local-cost}). Across the 60 measured calls there were
0 service exceptions, 0 timeouts and 0 calls with monitor errors.
The retained per-metric outcome table reports acceptance and rejection
alongside timing: syncopated-pedal rejected on both workloads, while the
other six metrics accepted on every measured call.
Separate original-monitor diagnostics locate these rejections at events
28 of 464 and
29 of 4946, respectively.
At adjacent pedal spans, a boundary attack is consumed by the preceding
region before its release; the following press lacks a fresh preceding attack.
These diagnostic calls are excluded from the timing samples.
Parsing warnings occurred on every call (1 occurrences per smaller-score
parse and 24 per larger-score parse), including ignored/unsupported
directions and an unmatched pedal stop; Prolog also emitted a singleton-variable
warning. All messages and raw samples are retained. An initial evaluation-wrapper
failure occurred before monitoring and is archived separately; it did not change
inputs, specifications or the frozen expectations.

\paragraph{Limits.}
The evidence concerns these synthetic cases, the original specifications and
two local workloads. A rejecting monitor may stop before consuming the whole
trace, so the results are not worst-case costs. They establish neither general
musical completeness nor independent alignment accuracy, scalability,
real-time guarantees or end-to-end user-facing latency. Scripts, MusicXML,
MIDI, service inputs, enriched traces, expected/actual verdicts, hashes,
warnings and individual timing samples accompany the evaluation.
